%======================================================================
\section{Beyond Branch-and-Price: A Taxonomy of IHS-OR Relationships}
\label{sec:taxonomy}
%======================================================================

While the branch-and-price and Benders/Dantzig--Wolfe analogies provide the
structural foundations for IHS, the algorithm can be viewed through several
other classical OR lenses. This section maps these relationships, with a focus
on the Lagrangian interpretation of the hitting-set search.

%----------------------------------------------------------------------
\subsection{The Lagrangian View: Discrete Dual Ascent}
\label{sec:lagrangian}
%----------------------------------------------------------------------

The most conceptual connection to Operations Research is the view of IHS as a
\textbf{discrete, logic-based form of Lagrangian Relaxation}.

In classical Lagrangian relaxation, one moves ``hard'' constraints into the
objective function, penalizing their violation with continuous multipliers
$\lambda$. For MaxSAT, the ``hard'' constraints are the SAT clauses
$\mathcal{H} \cup \mathcal{S}$. IHS relaxes the requirement to satisfy all soft
clauses and instead manages the penalties for falsifying them through the
hitting-set master.

\paragraph{IHS as a Lagrangian Dual solver.}
If we view the core-packing problem (\cref{eq:pack}) as the dual of the MaxSAT
relaxation, IHS is effectively solving the \emph{Lagrangian Dual} of the
hitting-set formulation. Each unsatisfiable core $\kappa$ represents a
combinatorial violation of the SAT constraints. In a continuous Lagrangian
setting (like subgradient descent), one would increase the multipliers $\lambda$
associated with these conflicts to penalize the violation. In IHS, we do not
adjust continuous multipliers; instead, we add a \emph{discrete hitting-set
constraint}. This constraint forces the master to ``pay'' for at least one
clause in the conflict, which is equivalent to a discrete step in a dual-ascent
process.

\paragraph{Closing the duality gap.}
A primary challenge in Lagrangian relaxation is the \emph{duality gap}: the
relaxed bound may never meet the integer optimum. IHS resolves this by
treating the SAT solver as a conflict generator that provides the necessary
\emph{binding} constraints to the dual. Because the hitting-set master is integer
(\textsc{IHS-IP}), it does not merely ascend to the LP bound $z^\star_{\mathrm{LP}}$
but continues adding cores until the lower bound meets an achievable upper bound,
effectively closing the integrality gap through iterative constraint generation.

\paragraph{Reduced costs and variable hardening.}
Practical solvers like MaxHS leverage this Lagrangian connection through
\emph{reduced-cost fixing}. By solving the LP relaxation of the current hitting
set master, the solver obtains dual values $y_i$ and reduced costs $\bar{c}_i$.
If $z^\star_{\mathrm{LP}} + \bar{c}_i > \mathit{ub}$ (the current best-known upper
bound), the soft clause $s_i$ \emph{must} be satisfied in any optimal solution.
This allows the solver to ``harden'' $s_i$ by moving it into $\mathcal{H}$, a
technique directly imported from the Lagrangian-based pruning used in MIP and
global optimization.

%----------------------------------------------------------------------
\subsection{The Cutting-Plane View: SAT as a Separation Oracle}
\label{sec:cutting}
%----------------------------------------------------------------------

From a primal perspective, IHS is a pure \textbf{Cutting Plane} algorithm. The
IP solver maintains a relaxed version of the problem (initially with no cores).
Each iteration generates a point $H^*$ that is ``optimal'' for the current
relaxation but ``infeasible'' for the full formula.

The SAT solver acts as the \textbf{Separation Oracle}. Given the current point
$H^*$, it searches for a linear inequality $\sum_{s_i \in \kappa} y_i \ge 1$
that is violated by $H^*$ (i.e., a core that is not hit). While classical OR
cutting planes (like Gomory cuts) are derived from the geometry of the LP
tableau, IHS cuts are \emph{combinatorial disjunctions} derived from the logical
structure of the SAT formula.

%----------------------------------------------------------------------
\subsection{The Dual Perspective: Column Generation in MSS Space}
\label{sec:mssspace}
%----------------------------------------------------------------------

The relationship between cores and hitting sets has a polar dual in the space of
\textbf{Maximal Satisfiable Subsets (MSSes)}. An assignment is optimal if it
satisfies an MSS of maximum total weight.

By the hitting-set duality of~\citet{liffiton2008algorithms}, a minimum hitting
set of all cores is equivalent to finding a Maximum Satisfiable Subset. Viewed
this way, IHS is a \textbf{Row Generation} algorithm on the hitting-set formulation.
Its dual is a \textbf{Column Generation} algorithm where each column represents an
MSS. This mirrors the Dantzig--Wolfe decomposition of \cref{sec:part1}, but where the
``columns'' are feasible assignments (MSSes) rather than dual variables for
cores.

%----------------------------------------------------------------------
\subsection{Summary Taxonomy}
\label{sec:taxonomy-summary}
%----------------------------------------------------------------------

\cref{tab:taxonomy} summarizes these multifaceted relationships, providing a
cross-community dictionary for the IHS approach.

\begin{table}[htbp]
\centering
\small
\caption{Taxonomy of IHS-OR Relationships.}
\label{tab:taxonomy}
\begin{tabular}{lll}
\toprule
\textbf{OR Perspective} & \textbf{Role of IHS Component} & \textbf{Key Mechanism} \\
\midrule
Logic-Based Benders & Master (IP) vs. Subproblem (SAT) & Logic-based Benders cuts \\
Cutting Planes      & SAT solver as Separation Oracle  & Combinatorial cuts \\
Dantzig--Wolfe      & Master as Dual Packing LP        & Column generation (cores) \\
Lagrangian Dual     & Master as Dual Optimizer         & Discrete dual ascent \\
Branch-and-Price    & Integer Master + Core Generation & Implicit branching \\
\bottomrule
\end{tabular}
\end{table}
